% !TEX root = ../bb_AiRa_Kappaleet.tex

%% HARRIS: Chapter 10

% ö = \%c3\%b6
% ä = \%C3\%A4

\xdef\sectiontitle{\IIsectiontitle} %aiheuttama
\TOCrefs

%%%%%%%%%%%%%%%
\begin{frame}[label=2_3_a]%[label=2_2_TOC1]
\frametitle{\texorpdfstring{\culine[\sectiontitleulinecolor]{\sectiontitle}}{\sectiontitle} }
\label{\TOCname}

\vspace{\chapterTOCvksip}
\begin{itemize}[leftmargin=0.5cm,itemsep=3mm]
    \item[2.1]<1-> \hyperlink{2_1_a}{\IIaaSubsectiontitle}
    \item[2.2]<1-> \hyperlink{2_2_a}{\IIaSubsectiontitle}
    \item[2.3]<1-> \hyperlink{2_3_a}{\CDAlert[red]{\IIbSubsectiontitle}}	
    \item[2.4]<1-> \hyperlink{2_4_a}{\IIcSubsectiontitle}	
    \item[2.5]<1-> \hyperlink{2_5_a}{\IIdSubsectiontitle}
    \item[2.6]<1-> \hyperlink{2_6_a}{\IIeSubsectiontitle}
    \item[2.7]<1-> \hyperlink{2_7_a}{\IIfSubsectiontitle}
    \item[2.8]<1-> \hyperlink{2_8_a}{\IIgSubsectiontitle}
    \item[2.9]<1-> \hyperlink{2_9_a}{\IIhSubsectiontitle}
    \item[2.10]<1-> \hyperlink{2_10_a}{\IIiSubsectiontitle}
\end{itemize}

%\endofslide 
\end{frame}
%%%%%%%%%%%%%%%

\xdef\sectiontitle{\IIbSubsectiontitle} 

%%%%%%%%%%%%%%%
\begin{frame}%[label=2_2_a]
\frametitle{\texorpdfstring{\culine[\sectiontitleulinecolor]{\sectiontitle}}{\sectiontitle} }
\framesubtitle{Absorptio- ja emissiospektrit}
\appear{}

\begin{itemize}
	\item Molekyylin elektroneilla on useita energiatasoja/-tiloja\appear{, minkä takia ne voivat absorboida ja emittoida sähkömagneettista säteilyä} 
	
	\item Molekyylien \CDAlert[red]{absorptio- ja emissiospektrit antavat tietoa} niiden \Alert[red]{energiatiloista} \appear{ja jopa \Alert[red]{rakenteesta}.}\appear{ Tätä tietoa voidaan hyödyntää arvioimaan esim. materiaalien \Alert[red]{lämpökapasiteetteja} }
\implies{ kemiallisilla yhdisteillä on \CDAlert[blue]{spektriset sormenjäljet} \appear{joiden avulla niitä voidaan tunnistaa ja mitata pitoisuuksia }\appear{(ts.~konsentraatioita)}}

\item Molekyylien energiaspektri sisältää tietoa:
\item[] \begin{itemize}
	\item Elektronitiloista ja -siirtymistä\appear{, ${\sim}$1–1000 eV} $\rightarrow$  \appear{${\sim}$\Alert[red]{näkyvä valo} ja \CDAlert[fuchsia]{röntgensäteet}}
\item Vibraatiotiloista\appear{, 1–100 meV} $\rightarrow$ \appear{mikroaallot ja \CDAlert[blue]{infrapuna}}
\item Rotaatiotiloista\appear{, 1–10 meV} $\rightarrow$ \appear{mikroaallot ja infrapuna}
\end{itemize}

\end{itemize}
%

\endofslide 
%\endofsection
\end{frame}
%%%%%%%%%%%%%%%

%%%%%%%%%%%%%%%
\begin{frame}
\frametitle{\texorpdfstring{\culine[\sectiontitleulinecolor]{\sectiontitle}}{\sectiontitle} }
\framesubtitle{Kaksiatominen molekyyli}
\appear{}

% 6.5 cm = midway, 10 cm = 3/4 
\vspace{1mm}
\begin{minipage}{9.2cm}
	\begin{itemize}
	\item Tarkastellaan \CDAlert[red]{kaksiatomista molekyyliä}.\\ \appear{Oletetaan potentiaalienergiakaivo likimain pa\-ra\-bo\-liseksi} \appear{tasapainopisteen $x_r$ läheisyydessä}
\[ 
 U \textrm{((}x_{r} \textrm{))} \appear{\cong} \appear{U_{0}}  \plus \frac{1}{2} \appear{\kappa x_r^{2}} \appear{\,,\qquad \text{missä}~x_{r}} \equiv \appear{x-a}
 \] 

\item Kerroin $\kappa$ \appear{on \CDAlert[tunired]{värähtelevän molekyylin} \\\CDAlert[red]{jousivakio} }
\[ 
 \kappa \equiv \frac{d^{2} U(x)}{d x^{2}} \appear{\,, \qquad kun~x}\approx \appear{a} 
 \]
\appear{Tällöin \CDAlert[red]{harmoniselle värähtelijälle} saadaan energiaksi} 
\[
E_{vib, n}  \equal   \appear{\textrm{((}n}  \plus  \frac{1}{2} \appear{\textrm{))}} \appear{\hbar \omega_{0}}
\]
\end{itemize}

\end{minipage}
\vspace{2.5mm}

\xdef\loc{north east}
\begin{tikzpicture}[remember picture,overlay] % anchor=south
    \node (A) [yshift=0.9*\figYshift,xshift=\figXshift, anchor=\loc] at (current page.\loc) {\includegraphics[page=1,width=4cm]{Figures_booklet/SOM_10_Bonding_figures/SOM_Bonding_Figure_16.png}}; %scale=\figscale. % 8cm max height
\end{tikzpicture}

\xdef\loc{south east}
\begin{tikzpicture}[remember picture,overlay] % anchor=south
    \node (A) [yshift=-0.35*\figYshift,xshift=\figXshift, anchor=\loc] at (current page.\loc) {\includegraphics[page=1,width=4cm]{Figures_booklet/SOM_10_Bonding_figures/SOM_Bonding_Figure_17.png}}; %scale=\figscale. % 8cm max height
\end{tikzpicture}
\footnote{\HarrisCRshort[1-]}

\endofslide 
%\endofsection
\end{frame}
%%%%%%%%%%%%%%%

%%%%%%%%%%%%%%%
\begin{frame}
\frametitle{\texorpdfstring{\culine[\sectiontitleulinecolor]{\sectiontitle}}{\sectiontitle} }
\framesubtitle{Kaksiatomisen molekyylin vibraatiotilat}
\appear{}

\begin{itemize}
	\item Mikä on kaksiatomisen molekyylin \CDAlert[red]{resonanssitaajuus $\omega_{0}$}? 
	
	\item värähtelytaajuus $\omega_{0}$ liittyy atomien välisiin voimiin. \appear{Jos kaksi massaa $m_{1}$} \appear{ja $m_{2}$} \appear{(}$\appear{\mu}\equiv$\appear{\,\CDAlert[red]{systeemin redusoitu massa})} \appear{on kytketty toisiinsa jousella}\appear{, systeemin potentiaalienergiaksi saadaan}
\[
\begin{aligned}
U(x)&  \equal \frac{1}{2} \appear{\kappa x^{2}}  \equal \frac{1}{2} \appear{\mu \omega_{0}{ }^{2} x^{2}} \\
\appear{\omega_{0}}&  \equal \appear{2 \pi f_{0}}  \equal \appear{\Sqrt{\Frac{\kappa}{\mu}} }\appear{\,, \qquad \mu}  \equal \fraca{m_{1} m_{2}}{m_{1}+m_{2}}
\end{aligned}
\]
\appear{ja \CDAlert[tunired]{vibraatiotiloiksi} muodostuu:}
\[ 
 E_{v i b, n}  \equal  \appear{\textrm{((}n}  \plus \frac{1}{2} \appear{\textrm{))}} \appear{\hbar} \appear{\Sqrt{\Frac{\kappa}{\mu}}} \appear{\,, \qquad n=0,1,2,3, \ldots }
 \]

\item Huomaa että tarkastelu on suoritettu yksiulotteiselle mallille
\end{itemize}

\endofslide 
%\endofsection
\end{frame}
%%%%%%%%%%%%%%%

%%%%%%%%%%%%%%%
\begin{frame}
\frametitle{\texorpdfstring{\culine[\sectiontitleulinecolor]{\sectiontitle}}{\sectiontitle} }
\framesubtitle{Kaksiatomisen molekyylin rotaatiotilat}
\appear{}

\seminormal
% 6.5 cm = midway, 10 cm = 3/4 
\vspace{1mm}
\begin{minipage}{9.5cm}
	\begin{itemize}[itemsep=1.6mm]
	\item Klassisesti pyörimiseen/\CDAlert[blue]{rotaatioon liittyvä kineettinen} \CDAlert[blue]{energia} on:\vspace{-3mm}  
%	$$
%	E=\frac{1}{2} I \omega^{2}=\frac{(I \omega)^{2}}{2 I}=\frac{L^{2}}{2 I} \,, \qquad where~I=moment~of~inertia
%	$$
%	 $$
%E=\frac{1}{2} I \omega^{2}=\frac{(I \omega)^{2}}{2 I}=\frac{L^{2}}{2 I} \,, \qquad 
%\begin{aligned}
% I&=moment~of~inertia  \\
% \omega&=angular~velocity~of~rotation \\
% L&=I \omega = molecule's~angular~momentum
%\end{aligned}
%$$ 
\[
\hspace{1.5cm} E  \equal \frac{1}{2} \appear{I} \appear{\omega^{2}}  \equal \frac{(I \omega)^{2}}{2 I}  \equal \frac{L^{2}}{2 I} \appear{\,,} 
\] \vspace{1mm} 
\appear{missä $I$ = hitausmomentti}\appear{, $\omega=$~kulmanopeus}\appear{, ja \\ $L=I \omega = $~kulmaliikemäärä}

\item Schrödingerin yhtälön ratkaisu pyörivälle kiin\-teäl\-le kappaleelle johtaa kulmaliikemäärän\\ \CDAlert[fuchsia]{kvantisoitumiseen} \appear{(kuten yksittäisillä atomeilla):}
\[ 
 L^{2}  \equal \appear{\ell} \appear{(\ell+1)} \appear{\hbar^{2}} \appear{\,, \qquad \ell=1,2,3, \ldots\Equiv \CDAlert[blue]{rotationaalinen kvanttiluku}} 
 \]
 \end{itemize}
\end{minipage}
%\vspace{2.5mm}

%Where $\ell$ is the rotational quantum number
\begin{itemize}
\item Huomaa että $L$ viittaa koko molekyyliin\appear{ joka pyörii massakeskipisteensä ympäri}\appear{, ja jolla on \CDAlert[tuniblue]{energiatasot}:}\vspace{-2mm} 
\[ 
 E_{r o t, \ell}  \equal \frac{\ell(\ell+1) \hbar^{2}}{2 I}  \equal \frac{\ell(\ell+1)}{2} \appear{E_{r}}  \quad  \qquad 
 \begin{aligned}
 \appear{E_r}  \equal  \frac{\hbar^2}{I}  \equal \appear{&~\text{karakteristinen}\\
 	&~\text{pyörimisenergia}}
 \end{aligned}
 \] %E_r = ~characteristic~rotational~energy of
\end{itemize}

\xdef\loc{north east}
\begin{tikzpicture}[remember picture,overlay] % anchor=south
    \node (A) [yshift=0.5*\figYshift,xshift=\figXshift, anchor=\loc] at (current page.\loc) {\includegraphics[page=1,width=4cm]{Figures_booklet/SOM_Tipler_Statistics_and_Bonding_figures/SOM_Tipler_figure_21.png}}; %scale=\figscale. % 8cm max height
\end{tikzpicture}
\footnoteleft{\TiplerCRshort[1-]}

\endofslide 
%\endofsection
\end{frame}
%%%%%%%%%%%%%%%

%%%%%%%%%%%%%%%
\begin{frame}
\frametitle{\texorpdfstring{\culine[\sectiontitleulinecolor]{\sectiontitle}}{\sectiontitle} }
\framesubtitle{Kaksiatomisen molekyylin rotaatiotilat}
\appear{}

\seminormal
\vspace{1mm}
\begin{minipage}{9.5cm}
	\begin{itemize}[itemsep=1.4mm]
	\item \CDAlert[blue]{Rotaatiotilat} sisältävät informaatiota mm. mo\-le\-kyy\-lien hitausmomentista \appear{eli myös \CDAlert[blue]{sidosten} \CDAlert[blue]{pituuksista}} 
	
	\item Pyörivän molekyylin hitausmomentti \appear{(massat $m_{1}$ ja $m_{2}$) on muotoa }
	\[ 
	I  \equal \appear{m_{1} r_{1}^{2}}  \plus \appear{m_{2} r_{2}^{2}}
	\]
\item \CDAlert[fuchsia]{Vipusäännön} (Arkhimedes!) \appear{$m_{1} r_{1}=m_{2} r_{2}$} \appear{ja $a \Equiv r_{1}+r_{2}$} \appear{avulla saadaan}
\[
I  \equal \appear{\mu a^{2}} \appear{\,,\qquad \text{missä redusoitu massa}~\mu}  \equal \frac{m_{1} m_{2}}{m_{1}+m_{2}}
\]
\end{itemize}
\vspace{1.25mm}
\end{minipage} 
\begin{itemize}[itemsep=1.6mm]
\item Yhtäsuurille massoille\appear{, kuten $\mathrm{H}_{2}$ tai $\mathrm{O}_{2}$}\appear{, saadaan $\mu=\Frac{1}{2} m_{1}  \equal \frac{1}{2} \appear{m_{2}}$} \appear{ja hitausmomentiksi tulee $I=\Frac{1}{2} m a^{2}$}

\item Yleisesti ottaen rotaatiotiloiksi saadaan \vspace{-0.7mm}
\[ 
 E  \equal \frac{\ell(\ell+1) \hbar^{2}}{2 I}  \equal \frac{\ell(\ell+1) \hbar^{2}}{2 \mu a^{2}} \equal \frac{\ell(\ell+1)}{2} \appear{E_{r}} \appear{\,\,, \qquad E_{r}=}\frac{\hbar^{2}}{\mu a^{2}} 
 \]
\end{itemize}

\xdef\loc{north east}
\begin{tikzpicture}[remember picture,overlay] % anchor=south
    \node (A) [yshift=0.5*\figYshift,xshift=\figXshift, anchor=\loc] at (current page.\loc) {\includegraphics[page=1,width=4cm]{Figures_booklet/SOM_Tipler_Statistics_and_Bonding_figures/SOM_Tipler_figure_21.png}}; %scale=\figscale. % 8cm max height
\end{tikzpicture}
\footnoteleft{\TiplerCRshort[1-]}

\endofslide 
%\endofsection
\end{frame}
%%%%%%%%%%%%%%%

%%%%%%%%%%%%%%%
\begin{frame}
\frametitle{\texorpdfstring{\culine[\sectiontitleulinecolor]{\sectiontitle}}{\sectiontitle} }
\framesubtitle{Kaksiatomisen molekyylin energiatilat}
\appear{}

\begin{itemize}[itemsep=1.6mm]
	\item \CDAlert[violet]{(Kaksiatomisen) molekyylin} energiatasot sisältävät siis tietoa \CDAlert[violet]{kahdesta kvanttiluvusta}: \appear{\CDAlert[red]{$n$}} \appear{ja \CDAlert[blue]{$\ell$}}\appear{, jotka liittyvät vastaavasti \CDAlert[red]{vibraatio-}} \appear{ja \CDAlert[blue]{rotaatio}energioihin:}
$$ \seminormal
\appear{E_{n, \ell}}  \equal \appear{E_{v i b}}  \plus \appear{E_{r o t}}  \equal  \appear{\textrm{((}n+\Frac{1}{2} \textrm{))}} \appear{\hbar} \appear{\Sqrt{\Frac{\kappa}{\mu}}} \plus \frac{\ell(\ell+1) \hbar^{2}}{2 \mu a^{2}}\appear{\,,} \quad \quad 
\begin{aligned}
&\appear{n=0,1,2,3, \ldots} \\
&\appear{\ell=0,1,2,3, \ldots}
\end{aligned}
$$
\vspace{-2mm}
\item Vaikka $\ell$ on seurausta samanlaisesta differentiaaliyhtälöstä \appear{kuin millä kuvataan vetyatomin elektronin kiertoliikettä}\appear{, se liittyy \CDAlert[blue]{molekyylin pyörimiseen}}

\item Kuten vedylle\appear{, erilaiset (kvantisoituneet) pyörimissuunnat ovat mahdollisia}\appear{, jota kuvataan kvanttiluvulla $m_{\ell}$} \appear{jolle pätee $m_{\ell}=0, \pm 1, \pm 2, \ldots, \pm \ell$} 

\item Huomaa myös että massa $\mu$ \appear{on molekyylin redusoitu massa} 

\item \CDAlert[violet]{Suurempien molekyylien} energiat\appear{ ja spektrit} \appear{\CDAlert[violet]{ovat vastaavanlaisia}}
\end{itemize}

\endofslide 
%\endofsection
\end{frame}
%%%%%%%%%%%%%%%

%%%%%%%%%%%%%%%
\begin{frame}
\frametitle{\texorpdfstring{\culine[\sectiontitleulinecolor]{\sectiontitle}}{\sectiontitle} }
\framesubtitle{Kaksiatomisen molekyylin energiatilat}
\appear{}

\begin{itemize}
	\item \CDAlert[blue]{Rotaatioenergiat}: 
	\begin{itemize} \small
	\item Elektronitilojen viritysenergiat ovat suuruusluokkaa $1$~eV\appear{, kun taas tyypillisten rotaatiotilojen energiat ovat luokkaa $10^{-4}-10^{-3}$~eV} 
	\item Transitiot rotaatiotilojen välillä synnyttävät fotoneita joiden aallonpituus on kaukoinfrapuna-alueella
	\item Rotaatioenergiat ovat myös pieniä \appear{verrattuna tyypilliseen} $\appear{k_{B} T} \approx \appear{0,026 ~\mathrm{eV}$} \appear{kun $T=300 \mathrm{~K}$} 
	\item \Alert[blue]{Normaaleissa lämpötiloissa molekyyli voi helposti virittyä alemmille rotaatiotiloille molekyylien törmätessä}

\end{itemize}

\item \CDAlert[red]{Vibraatioenergiat}: 
\begin{itemize} \small
	\item Normaaleissa lämpötiloissa \appear{vibraatioenergiat ovat tyypillisesti $0,1-0,2~\mathrm{eV}$ (eli $>k_{B} T $) }
	\item \Alert[red]{Suurin osa} molekyyleistä \appear{on siis \Alert[red]{alimmalla vibraatiotilalla}} \appear{missä $n=0$ ja $E_{0}=1 / 2 h f$}
\end{itemize}

\item \CDAlert[violet]{Huoneenlämpötilassa} \appear{molekyyli on elektronisessa perustilassa}\appear{, ja \CDAlert[red]{alimmalla  vibraatiotilalla}}\appear{, mutta voi olla \CDAlert[blue]{korkeammallakin rotaatiotilalla}}

\end{itemize}

\endofslide 
%\endofsection
\end{frame}
%%%%%%%%%%%%%%%

%%%%%%%%%%%%%%%
\begin{frame}
\frametitle{\texorpdfstring{\culine[\sectiontitleulinecolor]{\sectiontitle}}{\sectiontitle} }
\framesubtitle{Esimerkki}
\appear{}

%\vspace{2.5mm}
\begin{minipage}{5.75cm}
\begin{itemize}
\item \CDAlert[red]{Kysymys}: Tarkastellaan kaksiatomista HD molekyyliä \appear{(D = deuterium = vety + neutroni}\appear{, massat $m_H=1,007$ amu ja $m_D=2,013$ amu}\appear{, ja sidospituus on $0,074 \mathrm{~nm}$).} \appear{Käyttäen Boltzmannin jakaumaa}\appear{, määritä lämpötila jossa rotaatiotiloilla $\ell=1$ ja $\ell=0$} \appear{olevien tilojen populaatioiden suhde on $1 / 10$}
\end{itemize}
\end{minipage}
%\vspace{-2.5mm}
\begin{itemize}
\item[]
	
	\item \CDAlert[blue]{Vastaus}: \appear{Tilalla $\ell=1$ on kolme kvantittunutta rotaatiosuuntaa } \appear{($m_{\ell}$:n arvot)}\appear{, ja tilalla $\ell=0$ on vain yksi}\appear{. Populaatioiden suhteeksi saadaan siis}\vspace{-4mm}
\[ 
 \hspace{2cm} \frac{\text { Tilatiheys } E_{n=0, \ell=1}}{\text { Tilatiheys } E_{n=0, \ell=0}} \equal \frac{3}{1} \appear{\times} \frac{e^{-E_{0,1} / k_{B} T}}{e^{-E_{0,0} / k_{B} T}}  \equal \appear{3 e^{- \textrm{((}E_{0,1}-E_{0,0} \textrm{))} / k_{B} T} }
 \]
\end{itemize}

\xdef\loc{north east}
\begin{tikzpicture}[remember picture,overlay] % anchor=south
    \node (A) [yshift=0.9*\figYshift,xshift=\figXshift, anchor=\loc] at (current page.\loc) {\includegraphics[page=1,width=7.8cm]{Figures_booklet/SOM_10_Bonding_figures/SOM_Bonding_Figure_18.png}}; %scale=\figscale. % 8cm max height
\end{tikzpicture}
\footnote{\HarrisCR[1-]}

\endofslide 
%\endofsection
\end{frame}
%%%%%%%%%%%%%%%

%%%%%%%%%%%%%%%
\begin{frame}
\frametitle{\texorpdfstring{\culine[\sectiontitleulinecolor]{\sectiontitle}}{\sectiontitle} }
\framesubtitle{Esimerkki}
\appear{}

\begin{itemize}[itemsep=1.7mm]
	\item Molekyylin redusoitu massa on 
	$
\appear{\mu}  \equal \frac{1,007 \mathrm{\;amu} \,\times\, 2,013 \mathrm{\;amu}}{1,007 \mathrm{\;amu} \,+\,2,013 \mathrm{\;amu}}  \equal \appear{0,671 \mathrm{\;amu}}
$
	\appear{ja energia-ero on}
	$$
	 \begin{aligned}  \small
	 \appear{E_{0,1}-E_{0,0}} & \equal \frac{1(1+1) \hbar^{2}}{2 \mu a^{2}}  \minus \frac{0(0+1) \hbar^{2}}{2 \mu a^{2}}  \equal \frac{\hbar^{2}}{\mu a^{2}} \\
	 & \equal \fraca{ \textrm{((}1,055 \times 10^{-34} \,Js \textrm{))}^{2}}{0,671 \times 1,66 \times 10^{-27} \mathrm{\,kg} \times \textrm{((}0,074 \times 10^{-9} \mathrm{\,m} \textrm{))}^{2}}  \equal \appear{1,82 \times 10^{-21} J}  \equal \appear{0,011 \mathrm{\,eV}}
	 \end{aligned} 
	 $$ 
	
\item Asettamalla suhteeksi $1 / 10$ \appear{lämpötila voidaan ratkaista:}
\[ \small
\frac{1}{10}  \equal \appear{3 \exp}  \LB[3] \frac{-1,82 \times 10^{-21} \,J}{1,38 \times 10^{-23} \,J / K \times T} \RB[3] \qquad \Rightarrow \qquad \appear{ \cutline{red}{T \cong 40 \,K}}
\]
\vspace{-2mm}\implies{ \CDAlert[red]{Alle $40 \mathrm{~K}$, HD molekyylit voivat vain liikkua (eivät pyöriä!)} } 

\item Tämä käy yhteen kuvan 18 $\mathrm{H}_{2}$ molekyylien kanssa\appear{, joissa siirtymä translaatioenergioista rotaatioenergioihin tapahtuu vähän yli 40–50 K}
\end{itemize}

\endofslide 
%\endofsection
\end{frame}
%%%%%%%%%%%%%%%

%%%%%%%%%%%%%%%
\begin{frame}
\frametitle{\texorpdfstring{\culine[\sectiontitleulinecolor]{\sectiontitle}}{\sectiontitle} }
\framesubtitle{Valintasäännöt}
\appear{}

\begin{itemize}
	\item Vibraatiotilojen \appear{ja rotaatiotilojen \CDAlert[red]{energiat voivat saada vain} \CDAlert[red]{diskreettejä arvoja}}\appear{, minkä kuuluisi näkyä myös molekyylien absorptio- ja emissiospektreissä}

\item Jos fotoni emittoituu\appear{, osa molekyylin rotaatio- ja vibraatioenergiasta siirtyy fotonille}\appear{, ja $n$:n tai $\ell$:n tai molempien täytyy pienentyä}
\implies{ ~\CDAlert[red]{Valintasäännöt}: $\qquad \appear{\Delta \ell}  \equal \appear{\pm 1} \quad$ \appear{ja} $\appear{\quad \Delta n}  \equal \appear{\pm 1}$ }

\item  Fotonin spin on 1\appear{, joten systeemin kokonaiskulmaliikemäärä voi säilyä emissioprosessin aikana vain} \appear{jos $\ell$ muuttuu yhdellä}$\quad \rightarrow \quad \appear{\Delta \ell=\pm 1}$
\item Molekyylin varaustiheys värähtelee dipolin tavoin vibraatiotilojen välisissä sirtymissä. \appear{Varautunut harmoninen värähtelijä oskilloi} \appear{(ei osoiteta)} \appear{\CDAlert[blue]{sähködipolin tavoin} vain tilojen välisissä transitioissa joissa kvanttiluku muuttuu yhdellä} $\quad \rightarrow \quad \appear{\Delta n=\pm 1}$
\end{itemize}

\endofslide 
%\endofsection
\end{frame}
%%%%%%%%%%%%%%%

%%%%%%%%%%%%%%%
\begin{frame}
\frametitle{\texorpdfstring{\culine[\sectiontitleulinecolor]{\sectiontitle}}{\sectiontitle} }
\framesubtitle{Valintasäännöt}
\appear{}

\begin{itemize}
	\item Jos molekyylin vibraatioenergia pienenee\appear{, vibraatiotilan täytyy relaksoitua alemmalle elektroniselle tilalle ($\Delta n=-1$) }

\item Mutta koska vibraatiotilat ovat kauempana toisistaan kuin rotaatiotilat\appear{, vibraatiotilan relaksoituminen}\appear{ voi johtaa joko rotaatiotilan kasvuun tai pienenemiseen.}\appear{ Toisin sanoen \CDAlert[red]{$\ell$ voi joko pienentyä tai kasvaa yhdellä} }

\item Emittoituneen fotonin energia on siis yhteensä muotoa
\[ 
 E_{\text {fotoni}}  \equal \appear{\hbar} \appear{\Sqrt{\Frac{\kappa}{\mu}}} \appear{\pm} \appear{i \times} \frac{\hbar^{2}}{\mu a^{2}}  \appear{\,, \qquad i=1,2,3, \ldots }
 \]
\end{itemize}

\endofslide 
%\endofsection
\end{frame}
%%%%%%%%%%%%%%%

%%%%%%%%%%%%%%%
\begin{frame}
\frametitle{\texorpdfstring{\culine[\sectiontitleulinecolor]{\sectiontitle}}{\sectiontitle} }
\framesubtitle{Valintasäännöt}
\appear{}

\begin{itemize}[itemsep=1.8mm]
	\item Kuvassa 19\appear{ näkyy vibraatiotilat $n=0$ ja $n=1$}\appear{, sekä rotaatiotilat $\ell=0,1,2$ ja 3}\appear{, siis $\Frac{\ell(\ell+1)}{2}=0,1,3$ ja 6} 

\item Me \Alert[red]{löydämme kuusi transitiota jotka toteuttavat valintasäännöt} 
\end{itemize}

% 6.5 cm = midway, 10 cm = 3/4 
\vspace{1mm}
\begin{minipage}{7cm}
\begin{itemize}[itemsep=1.8mm]
	\item Huomaa että rotaatiotilojen välisen energia-eron kvantti on
	\[
	\delta E_{r}  \equal \frac{\hbar^{2}}{\mu a^{2}}
	\]

\item Spektrin keskellä on 'aukko'\appear{, koska transitiota $\Delta \ell=0$ ei tapahdu}\appear{, se on kielletty}

\item Säännöt, joita nuolet kuvassa selventävät\appear{, pätevät \CDAlert[red]{sekä } \CDAlert[red]{absorptiolle}} \appear{että \CDAlert[red]{emissiolle}}
\end{itemize}
\end{minipage}


\xdef\loc{south east}
\begin{tikzpicture}[remember picture,overlay] % anchor=south
    \node (A) [yshift=-0.25*\figYshift,xshift=\figXshift, anchor=\loc] at (current page.\loc) {\includegraphics[page=1,height=6.4cm]{Figures_booklet/SOM_10_Bonding_figures/SOM_Bonding_Figure_19.png}}; %scale=\figscale. % 8cm max height
\end{tikzpicture}
\footnote{\HarrisCR[1-]}

\endofslide 
%\endofsection
\end{frame}
%%%%%%%%%%%%%%%

%%%%%%%%%%%%%%%
\begin{frame}
\frametitle{\texorpdfstring{\culine[\sectiontitleulinecolor]{\sectiontitle}}{\sectiontitle} }
\framesubtitle{Rotaatio–vibraatiospektri}
\appear{}

\begin{itemize}[itemsep=1.8mm]
	\item Mahdollisten transitioiden joukko\appear{, tai vastaavasti emittoituneiden fotonien energiat} \appear{(tai absorboituneiden fotonien)}\appear{, muodostavat \CDAlert[red]{rotaatio–vibraatiospektrin}}
	\item Transitioita joissa sekä vibraatiotila \appear{että rotaatiotila muuttuu} \appear{kutsutaan  \CDAlert[red]{rotaatio–vibraatiotransitioksi}}
\appear{\xdef\loc{south east}
\begin{tikzpicture}[remember picture,overlay] % anchor=south
    \node (A) [yshift=-0.35*\figYshift,xshift=\figXshift, anchor=\loc] at (current page.\loc) {\includegraphics[page=1,width=7.5cm]{Figures_booklet/SOM_10_Bonding_figures/SOM_Bonding_Figure_20.png}}; %scale=\figscale. % 8cm max height
\end{tikzpicture}\footnote{\HarrisCR[1-]}}
\item Kuvassa~20 näemme $\mathrm{HCl}$:n absorptiospektrin\appear{, jossa absorptioviivat } \appear{ muodostavat ylläkuvatun rakenteen.}\appear{ Jopa spektrin keskellä olevan \CDAlert[blue]{'aukon'}}
\end{itemize}

% 6.5 cm = midway, 10 cm = 3/4
\vspace{2mm} 
\begin{minipage}{6.8cm}
\begin{itemize}[itemsep=1.8mm]
	\item[] \begin{itemize} \small
	 \item \CDAlert[tunired]{Piikit eivät ole yhtä kaukana toisistaan}\appear{, koska $\ell$:n kasvaessa}\appear{ molekyylit venyvät } \appear{muuttaen niiden energioita}
	\item Piikit ovat \CDAlert[tunired]{leventyneet}\appear{, koska kloorilla on kaksi yleistä  \CDAlert[fuchsia]{isotooppia}} \appear{${ }^{35} \mathrm{Cl}$ ja ${ }^{37} \mathrm{Cl}$} \appear{joilla on poikkeavat massat } \appear{(muista: $M_{C l}=35,45 \mathrm{~g} / \mathrm{mol} $)}
\end{itemize}
\end{itemize}
\end{minipage}

\endofslide 
%\endofsection
\end{frame}
%%%%%%%%%%%%%%%

%% LET's SKIP THE REAL TRANSMISSION SPECTRUM OF HCl SLIDE for time being
%%%%%%%%%%%%%%%%
%\begin{frame}
%\frametitle{\texorpdfstring{\culine[\sectiontitleulinecolor]{\sectiontitle}}{\sectiontitle} }
%
%\begin{itemize}
%	\item Note, that the absorption spectrum of HCl shown in Fig.~20 was actually processed from the real transmission spectrum, by first removing the background and further inverting the $y$-axis
%\end{itemize}
%
%\xdef\loc{east}
%\begin{tikzpicture}[remember picture,overlay] % anchor=south
%    \node (A) [yshift=\figYshift,xshift=\figXshift, anchor=\loc] at (current page.\loc) {\includegraphics[page=1,width=6cm]{Figures_booklet/SOM_10_Bonding_figures/SOM_Bonding_Figure_20.png}}; %scale=\figscale. % 8cm max height
%\end{tikzpicture}
%
%\xdef\loc{west}
%\begin{tikzpicture}[remember picture,overlay] % anchor=south
%    \node (A) [yshift=0*\figYshift,xshift=-\figXshift, anchor=\loc] at (current page.\loc) {\includegraphics[page=1,width=6cm]{Figures_booklet/SOM_10_Bonding_figures/SOM_Bonding_Figure_20.png}}; %scale=\figscale. % 8cm max height
%\end{tikzpicture}
%
%\xdef\loc{east}
%\begin{tikzpicture}[remember picture,overlay] % anchor=south
%    \node (A) [yshift=\figYshift,xshift=\figXshift, anchor=\loc] at (current page.\loc) {\includegraphics[page=1,width=6cm]{Figures_booklet/SOM_10_Bonding_figures/SOM_Bonding_Figure_20.png}}; %scale=\figscale. % 8cm max height
%\end{tikzpicture}
%
%
%\endofslide 
%%\endofsection
%\end{frame}
%%%%%%%%%%%%%%%%

%%%%%%%%%%%%%%%
\begin{frame}
\frametitle{\texorpdfstring{\culine[\sectiontitleulinecolor]{\sectiontitle}}{\sectiontitle} }
\framesubtitle{Energioiden suuruusluokista}
\appear{}

\begin{itemize}[itemsep=1.6mm]
	\item Elektronitilojen väliset energiaerot ovat molekyyleille usein $>\!1~\mathrm{eV}$\appear{, vibraatiotiloille $>\!0,1~\mathrm{eV}$} \appear{ja rotaatiotiloille $<\!0,1~ \mathrm{eV}$ }
	
	\item Kuva 21 ei ole mittakaavassa\appear{, mutta kuvaa hyvin tilojen energioiden välisiä suuruusluokkia}
	
% 6.5 cm = midway, 10 cm = 3/4 
%\vspace{2.5mm}
\item[] \begin{minipage}{8.2cm} \small
 	\begin{itemize}
	\item  Punaiset käyrät esittävät kahta elektronista tilaa
	
	\item Kun elektroni on \CDAlert[fuchsia]{virittyneessä tilassa}\appear{, atomit ovat kauempana toisistaan} \appear{($a$ muuttuu)} \appear{mutta silti sidoksissa toisiinsa} 
	
	\item Potentiaalikaivo on leveämpi ja matalampi \appear{(= kuten kuvan ylempi kaivo)}
	
	\item \CDAlert[red]{Kummallakin elektronitilalla}\appear{ on useita vib\-raa\-tiotiloja $n$} \appear{ja vielä useampia ja lähempänä toisiaan olevia rotaatiotiloja $\ell$}
\end{itemize}

\end{minipage}
\vspace{2.5mm}

\end{itemize}

\xdef\loc{south east}
\begin{tikzpicture}[remember picture,overlay] % anchor=south
    \node (A) [yshift=-0.35*\figYshift,xshift=\figXshift, anchor=\loc] at (current page.\loc) {\includegraphics[page=1,height=6.1cm]{Figures_booklet/SOM_10_Bonding_figures/SOM_Bonding_Figure_21.png}}; %scale=\figscale. % 8cm max height
\end{tikzpicture}
\footnote{\HarrisCR[1-]}

\endofslide 
%\endofsection
\end{frame}
%%%%%%%%%%%%%%%

%%%%%%%%%%%%%%%
\begin{frame}
\frametitle{\texorpdfstring{\culine[\sectiontitleulinecolor]{\sectiontitle}}{\sectiontitle} }
\framesubtitle{Monimutkaisemmat molekyylit}
\appear{}

\begin{itemize}
	\item Huomaa että olemme tarkastelleet vain \CDAlert[red]{kaksiatomista molekyyliä} 
	
	\item Suurimmalla osalla molekyyleistä \CDAlert[blue]{on monimutkaisempi rakenne}\appear{, ja vastaavasti sekä niiden rotaatio-} \appear{että vibraatioenergiaspektrit \CDAlert[tuniblue]{ovat monimutkaisempia}}

\item Esimerkiksi $\mathrm{CO}_{2}$ (vasen) \appear{ja $\mathrm{H}_{2} \mathrm{O}$ (oikea)} \appear{ovat kaksi yleistä molekyyliä} \appear{jotka voivat värähdellä monella eri tavalla, kuten kuvassa näkyy}

\end{itemize}

\appear{\xdef\loc{south}
\begin{tikzpicture}[remember picture,overlay] % anchor=south
    \node (A) [yshift=-0.5*\figYshift,xshift=0*\figXshift, anchor=\loc] at (current page.\loc) {\includegraphics[page=1, height=4cm]{Figures_booklet/SOM_10_Bonding_figures/Tikz_SOM_Bonding_vibrational_modes_fixed}};
\end{tikzpicture}}

%\endofslide 
\endofsection
\end{frame}
%%%%%%%%%%%%%%%